% 	Name		:: 	sthlm Beamer Theme  HEAVILY based on the hsrmbeamer theme (Benjamin Weiss)
%	Author		:: 	Mark Hendry Olson (mark@hendryolson.com)
%	Created		::	2013-07-31
%	Updated		::	June 18, 2015 at 08:45
%	Version		:: 	1.0.2
%	Email		:: 	hendryolson@gmail.com
%	Website		:: 	http://v42.com
%
% 	License		:: 	This file may be distributed and/or modified under the
%                  	GNU Public License.
%
%	Description	::	This presentation is a demonstration of the sthlm beamer
%					theme, which is HEAVILY based on the HSRM beamer theme created by Benjamin Weiss
%					(benjamin.weiss@student.hs-rm.de), which can be found on GitHub
%					<https://github.com/hsrmbeamertheme/hsrmbeamertheme>.


%-=-=-=-=-=-=-=-=-=-=-=-=-=-=-=-=-=-=-=-=-=-=-=-=
%
%        LOADING DOCUMENT
%
%-=-=-=-=-=-=-=-=-=-=-=-=-=-=-=-=-=-=-=-=-=-=-=-=

\documentclass[newPxFont]{beamer}
\usetheme{sthlm}
%\usecolortheme{sthlmv42}

%-=-=-=-=-=-=-=-=-=-=-=-=-=-=-=-=-=-=-=-=-=-=-=-=
%        LOADING PACKAGES
%-=-=-=-=-=-=-=-=-=-=-=-=-=-=-=-=-=-=-=-=-=-=-=-=
\usepackage[utf8]{inputenc}
\usepackage[T1]{fontenc}
\usepackage{subfigure}
\usepackage{verbatim}

\newcommand{\tabitem}{%
  \usebeamertemplate{itemize item}\hspace*{\labelsep}}

%\renewcommand{\event}[3][e]{%
%  \pgfmathsetlength\xstop{(#2-\theyearstart)*\unit}%
%  \ifx #1e%
%    \draw[fill=black,draw=none,opacity=0.5]%
%      (\xstop, 0) circle (.2\unit)%
%      node[opacity=1,rotate=45,right=.2\unit] {#3};%
%  \else%
%    \pgfmathsetlength\xstart{(#1-\theyearstart)*\unit}%
%    \draw[fill=black,draw=none,opacity=0.5,rounded corners=.1\unit]%
%      (\xstart,-.1\unit) rectangle%
%      node[opacity=1,rotate=45,right=.2\unit] {#3} (\xstop,.1\unit);%
%  \fi}%

%-=-=-=-=-=-=-=-=-=-=-=-=-=-=-=-=-=-=-=-=-=-=-=-=
%        BEAMER OPTIONS
%-=-=-=-=-=-=-=-=-=-=-=-=-=-=-=-=-=-=-=-=-=-=-=-=

%\setbeameroption{show notes}

%-=-=-=-=-=-=-=-=-=-=-=-=-=-=-=-=-=-=-=-=-=-=-=-=
%
%	PRESENTATION INFORMATION
%
%-=-=-=-=-=-=-=-=-=-=-=-=-=-=-=-=-=-=-=-=-=-=-=-=

\title{Un systeme de base de donnees pour la Grande Muraille Verte}
\subtitle{Un outil pour la recherche et la gestion}
%\date{\small{\jobname}}
%\date{\today}
\date{26 octobre 2017}
\author{\texttt{ E. Delay}}
\institute{\textsc{Ohm} Téssékéré - Sénégal}

\hypersetup{
pdfauthor = {E. DELAY},
pdfsubject = {Réseau OHM},
pdfkeywords = {grande muraille verte, ANR future Sahel},
pdfmoddate= {D:\pdfdate},
pdfcreator = {}
}

\begin{document}

%-=-=-=-=-=-=-=-=-=-=-=-=-=-=-=-=-=-=-=-=-=-=-=-=
%
%	TITLE PAGE
%
%-=-=-=-=-=-=-=-=-=-=-=-=-=-=-=-=-=-=-=-=-=-=-=-=

\maketitle

%\begin{frame}[plain]
%	\titlepage
%\end{frame}

%-=-=-=-=-=-=-=-=-=-=-=-=-=-=-=-=-=-=-=-=-=-=-=-=
%
%	TABLE OF CONTENTS: OVERVIEW
%
%-=-=-=-=-=-=-=-=-=-=-=-=-=-=-=-=-=-=-=-=-=-=-=-=
%\section*{Overview}
%\begin{frame}{Overview}
%% For longer presentations use hideallsubsections option
%\tableofcontents[hideallsubsections]
%\end{frame}

%-=-=-=-=-=-=-=-=-=-=-=-=-=-=-=-=-=-=-=-=-=-=-=-=
%	FRAME: INTRODUCTION
%-=-=-=-=-=-=-=-=-=-=-=-=-=-=-=-=-=-=-=-=-=-=-=-=

\section{Introduction}

%-=-=-=-=-=-=-=-=-=-=-=-=-=-=-=-=-=-=-=-=-=-=-=-=
%	FRAME: ANR framework in 4 wp
%-=-=-=-=-=-=-=-=-=-=-=-=-=-=-=-=-=-=-=-=-=-=-=-=
\begin{frame}[c]{La grande Muraille verte}
\vspace{-1cm}
Juste pour se remettre dans le bain ! L'ANR \textit{Future Sahel} se focalise sur des  approches multi-scalaires pour de meilleures pratiques de gestion des ressources des espaces sahéliens de la Grande Muraille Verte.
\begin{figure}
	\centering
	\includegraphics[width = 0.5\textwidth]{img/great_wall}
  \includegraphics[width = 0.5\textwidth]{img/GMV_widou}
\end{figure}
\begin{center}
    \begin{tabular}{@{}l@{}}
       \\
      \tabitem width $\approx 15 km$ \\
      \tabitem length $\approx 7600km$
  \end{tabular}
\end{center}
\end{frame}

%-=-=-=-=-=-=-=-=-=-=-=-=-=-=-=-=-=-=-=-=-=-=-=-=
%	FRAME: ANR framework in 4 wp
%-=-=-=-=-=-=-=-=-=-=-=-=-=-=-=-=-=-=-=-=-=-=-=-=
\begin{frame}[c]{Future Sahel Framework}
\vspace{-1cm}
\begin{figure}
	\centering
	\includegraphics[width = 0.8\textwidth]{img/Carte_FutureSahel}
	\caption{Les Works packages de \textit{Future Sahel}}
\end{figure}
\end{frame}

%-=-=-=-=-=-=-=-=-=-=-=-=-=-=-=-=-=-=-=-=-=-=-=-=
%	FRAME: ANR framework in 4 wp
%-=-=-=-=-=-=-=-=-=-=-=-=-=-=-=-=-=-=-=-=-=-=-=-=

\begin{frame}[c]{Objectifs}
\vspace{-1cm}
En coopération étroite avec l'agence de la Grande Muraille Verte Sénégalaise, l'objectif est de coconstruire une base de données à destination de la recherche et du développement.
	\begin{itemize}
		\item Les chercheurs $\rightarrow$ maintiennent et exploitent des données produites dans un contexte de recherche;
		\item Les gestionnaires $\rightarrow$ utilisent leurs connaissances pour de la planification spatiale et temporelle.
	\end{itemize}
  \begin{figure}
  	\centering
  	\includegraphics[width = 6cm]{img/IR_spot}
  \end{figure}
\end{frame}

%-=-=-=-=-=-=-=-=-=-=-=-=-=-=-=-=-=-=-=-=-=-=-=-=
%	FRAME: Goals for our SGBD
%-=-=-=-=-=-=-=-=-=-=-=-=-=-=-=-=-=-=-=-=-=-=-=-=
\begin{frame}[c]{Vous avez dit base de données?}
\vspace{-1cm}
Est-ce qu'on peut rendre une présentation de base de données attrayante ?
\begin{figure}
	\centering
	\includegraphics[height=4cm]{img/sexyadmin}~
  \includegraphics[height=4cm]{img/lordofthering}
\end{figure}
\end{frame}

\begin{frame}[c]{Materiel et méthodes}
\vspace{-1cm}
\begin{itemize}
	\item L'architecture est basée sur PostgreSQL et PostGIS (GEOLAB);
	\item \^Etre compatible avec BBees metadata;
	\item Besoin d'être compatible entre données hétérogènes spatiales et temporelle :
	\begin{itemize}
		\item Raster (Spot, Modis, Landsat, Sentinel);
		\item Des données statistiques produites par les institutions;
		\item Données de terrain.
	\end{itemize}
	\item Les technologies doivent être transférées aux acteurs (ANGMV). Nous devons choisir des logiciels libres et \textit{open source}.
\end{itemize}
\begin{figure}
	\centering
	\includegraphics[height=20mm]{img/PostGIS_logo}
	\includegraphics[height=20mm]{img/QGis_Logo}
	\includegraphics[height=20mm]{img/GrassGIS_banner}
	\includegraphics[height=20mm]{img/Rlogo}
\end{figure}
\end{frame}

\begin{frame}[c]{Une architecture mondiale}
\vspace{-1cm}
\begin{itemize}
	\item La base de données est hébergé a deux endroits :
	\begin{enumerate}
		\item Dakar - UMI ;
		\item Clermont-Ferrand - Centre Régional de Ressources informatiques ;
	\end{enumerate}
	\item Dans le futur, une troisième version de la base sera hébergé à l'ANGMV.
\end{itemize}
\begin{figure}
	\centering
	\includegraphics[height=5cm]{img/sch_database}
\end{figure}
\end{frame}

%-=-=-=-=-=-=-=-=-=-=-=-=-=-=-=-=-=-=-=-=-=-=-=-=
%	FRAME: context
%-=-=-=-=-=-=-=-=-=-=-=-=-=-=-=-=-=-=-=-=-=-=-=-=

\section{Pathway}

%-=-=-=-=-=-=-=-=-=-=-=-=-=-=-=-=-=-=-=-=-=-=-=-=
%	FRAME: WP1 GENERAL
%-=-=-=-=-=-=-=-=-=-=-=-=-=-=-=-=-=-=-=-=-=-=-=-=

\begin{frame}[c]{WP1 : Des données "généralistes"}
\vspace{-2em}
Lister les données en utilisant le framework \textit{Institutional analysis and development} (IDA) (Ostrom 2009)
\vspace{-1em}
\begin{figure}
	\centering
	\includegraphics[width = 7cm]{img/IDA_data_org}
\end{figure}
\end{frame}

%-=-=-=-=-=-=-=-=-=-=-=-=-=-=-=-=-=-=-=-=-=-=-=-=
%	FRAME: WP1 Biophy
%-=-=-=-=-=-=-=-=-=-=-=-=-=-=-=-=-=-=-=-=-=-=-=-=

\begin{frame}[c]{WP1 : Données Biophysique}
\vspace{-2em}
\begin{itemize}
  \item Work in small windows
  \begin{itemize}
    \item Spot 6 images (1.5 m) $\rightarrow$ détection de la canopé, des mares, et NDVI;
    \item Sentinel 2 (10 m) $\rightarrow$ unités paysagères;
    \item Landsat (30 m) $\rightarrow$ unités paysagères;
    \item MODIS (250 m) $\rightarrow$ évaluation de la participation des arbres au NDVI;
    \item Grimms (5km) $\rightarrow$ reverdissement ?
  \end{itemize}
  \item Generalization $\rightarrow$ Sentinel ?
\end{itemize}
\vspace{-1em}
\begin{figure}
	\centering
	\includegraphics[width = 5cm]{img/NDVI2015}
  \includegraphics[width = 5cm]{img/reverdiement}
\end{figure}
\end{frame}

\begin{frame}[c]{WP1 : Calcul à partir de données publiques}
\vspace{-2em}
Usage abondant de \textit{google earth engine} pour effectuer des calculs sur des gros jeux de données publiques $\rightarrow$ intégrer ces données dans la base de données.
\vspace{-1em}
\begin{figure}
	\centering
	\includegraphics[width = 10cm]{img/ggearthengine}
\end{figure}
\end{frame}
%-=-=-=-=-=-=-=-=-=-=-=-=-=-=-=-=-=-=-=-=-=-=-=-=
%	FRAME: WP
%-=-=-=-=-=-=-=-=-=-=-=-=-=-=-=-=-=-=-=-=-=-=-=-=

\begin{frame}[c]{WP2 : Biodiversité cultivée}
\vspace{-2em}
Maximiser la biodiversité plantée en utilisant le protocole de reforestation de la GMV. Initialement les plantations n'étaient pas très diversifiées, principalement composées des espèces d'arbres: \textit{Balanites aegyptiaca}, \textit{Acacia senegal}, and \textit{Acacia tortilis subsp raddiana}.
\begin{itemize}
  \item Evaluer la plantation de nouvelles parcelles de reboisement,
  \item Monitoring par télédétection des essais de plantation (\textit{in situ}).
\end{itemize}
\begin{figure}
	\centering
	\includegraphics[width = 4cm]{img/Trees_sup15m2}
  \includegraphics[width = 4cm]{img/GRID_Trees_sup15m2}
\end{figure}
\end{frame}
%-=-=-=-=-=-=-=-=-=-=-=-=-=-=-=-=-=-=-=-=-=-=-=-=
%	FRAME: WP
%-=-=-=-=-=-=-=-=-=-=-=-=-=-=-=-=-=-=-=-=-=-=-=-=

\begin{frame}[c]{WP3 : Valorisation de la "filère" \textit{B. aegyptiaca}}
\vspace{-2em}
Cette espèce est utilisée par les acteurs pour  : les fruits, l'huile d'amande (cuisine et cosmétique).
\begin{itemize}
  \item Etude biogéographique,
  \item phénologie et production (biomasse, fruits),
  \item variabilité intraspécifique
\end{itemize}
\vspace{-1em}
\begin{figure}
	\centering
	\includegraphics[width = 7cm]{img/Biomasse_arboree}
  \vspace{-1em}
  \caption{Biomasse arborée (J.-L.Peiry and D. Ndiaye)}
\end{figure}
\end{frame}

\begin{frame}[c]{WP4 : Résilience et gestion "territoriale"}
\vspace{-2em}
Construire une cadre conceptuel de gestion des ressources naturelles basé sur l'approche "résilience" socio-environnementale et services écosystémiques (SES).
\begin{itemize}
  \item Analyse historique des projets de développement (Mazzero 2017)
  \item Cartographie des SES le long de la GMV,
  \item Améliorer la résilience du système.
\end{itemize}
\begin{figure}
	\centering
	\includegraphics[height = 3.5cm]{img/ggearth}
  \includegraphics[height = 3.5cm]{img/Sinare_et_al2016}
  \caption{\small{Screenshot Google-earth and maps from Sinare et al. (2016)}}
\end{figure}
\end{frame}

%-=-=-=-=-=-=-=-=-=-=-=-=-=-=-=-=-=-=-=-=-=-=-=-=
%	FRAME: LA BASE DE DONNEES AUJOURD'HUI
%-=-=-=-=-=-=-=-=-=-=-=-=-=-=-=-=-=-=-=-=-=-=-=-=
\section{La base de donnees aujourd'hui}

%-=-=-=-=-=-=-=-=-=-=-=-=-=-=-=-=-=-=-=-=-=-=-=-=
%	FRAME: UN OUTIL D'EMPOWERMENT
%-=-=-=-=-=-=-=-=-=-=-=-=-=-=-=-=-=-=-=-=-=-=-=-=
\begin{frame}[c]{Un outil d'\textit{empowerment}!}
\vspace{-2em}
\begin{figure}
	\centering
	\includegraphics[height = 4.5cm]{img/gandalf}
\end{figure}
\end{frame}

%-=-=-=-=-=-=-=-=-=-=-=-=-=-=-=-=-=-=-=-=-=-=-=-=
%	FRAME: EN LIGNE
%-=-=-=-=-=-=-=-=-=-=-=-=-=-=-=-=-=-=-=-=-=-=-=-=
\begin{frame}[c]{Architecture client/serveur}
\vspace{-2em}
\begin{figure}
	\centering
	\includegraphics[width = 0.9\textwidth]{img/cooperation}
\end{figure}
\end{frame}

%-=-=-=-=-=-=-=-=-=-=-=-=-=-=-=-=-=-=-=-=-=-=-=-=
%	FRAME: La base de données en Chiffre
%-=-=-=-=-=-=-=-=-=-=-=-=-=-=-=-=-=-=-=-=-=-=-=-=
\begin{frame}[fragile]{La base de données en chiffre}
\vspace{-2em}
\begin{itemize}
  \item Pas d'interface web !
  \item 85 tables spatailes provenant de :
    \begin{itemize}
      \item Agence Nationnal d'Amenagement du territoire,
      \item CSE,
      \item Prodig,
      \item FAO,
      \item NASA, Copernicus, etc,
      \item Données de recherches.
    \end{itemize}
\end{itemize}
\begin{figure}
	\centering
	\includegraphics[width = 0.9\textwidth]{img/data_golds}
\end{figure}
\end{frame}

%-=-=-=-=-=-=-=-=-=-=-=-=-=-=-=-=-=-=-=-=-=-=-=-=
%	FRAME: MANIPULATION
%-=-=-=-=-=-=-=-=-=-=-=-=-=-=-=-=-=-=-=-=-=-=-=-=
\begin{frame}[fragile]{Connecter R et postgreSQL}
\vspace{-2em}
\begin{verbatim}
require(RPostgreSQL)
m <- dbDriver("PostgreSQL")
con <- dbConnect(m, dbname="fs_gis",
          host='172.17.0.2',port='5432',
          user='postgres',password='postgres')
q <- "SELECT ST_AsText(geom) AS geom
      FROM data_gmv.communes_gmv
      LIMIT 10;"
rs <- dbSendQuery(con,q)
df <- fetch(rs,n=-1)
\end{verbatim}
\end{frame}

\begin{frame}[fragile]{Connecter R et postGIS}
\vspace{-2em}
\begin{verbatim}
library(rpostgis)
con<-dbConnect("PostgreSQL",
                dbname='fs_gis',
                host='172.17.0.2',
                port='5432',
                user='postgres',
                password='postgres')
shp <- pgGetGeom(con,
            name=c("data_gmv","communes_gmv"),
            geom = "geom")
plot(shp)
\end{verbatim}
\end{frame}


%-=-=-=-=-=-=-=-=-=-=-=-=-=-=-=-=-=-=-=-=-=-=-=-=
%	FRAME: connexion
%-=-=-=-=-=-=-=-=-=-=-=-=-=-=-=-=-=-=-=-=-=-=-=-=
\begin{frame}[c]{Se connecter avec Qgis}
\vspace{-2em}
\begin{figure}
	\centering
	\includegraphics[width = 0.8\textwidth]{img/qgis_screenshot}
\end{figure}
\end{frame}

%-=-=-=-=-=-=-=-=-=-=-=-=-=-=-=-=-=-=-=-=-=-=-=-=
%	FRAME: METADONNEES
%-=-=-=-=-=-=-=-=-=-=-=-=-=-=-=-=-=-=-=-=-=-=-=-=
\begin{frame}[c]{Un serveur de métadonnées}
\vspace{-2em}
\begin{figure}
	\centering
	\includegraphics[width = 0.6\textwidth]{img/geonetwork}
\end{figure}
Le travail de création de métadonnées est à poursuivre (reste 70\%).
\end{frame}

%-=-=-=-=-=-=-=-=-=-=-=-=-=-=-=-=-=-=-=-=-=-=-=-=
%	FRAME: Conclusion
%-=-=-=-=-=-=-=-=-=-=-=-=-=-=-=-=-=-=-=-=-=-=-=-=
\section{Perspective}

\begin{frame}[c]{Database modelling and ComMod}
\vspace{-2em}
Créer une base de données est un processus itératif. Dans le cadre de \textit{Future Sahel}, cette construction est en fait une co-construction avec les acteurs ... ça prend du temps !
\begin{figure}
	\centering
	\includegraphics[height = 3cm]{img/iteration}
\end{figure}
La démarche \includegraphics[height = 0.4cm]{img/logo_commod} : \textbf{i)} vocation à être remises en cause, \textbf{ii)} aucune hypothèse implicite, \textbf{iii)} processus de
validation.
\end{frame}

\begin{frame}[c]{Quid des données spécifique de recherche ?}
\vspace{-2em}
\begin{block}{\textsc{Comment aller plus loin ?}}
  Est-ce qu'il est possible de construire une ontologie commune aux chercheurs et au gestionnaire, pour proposer un schéma de base donnée suffisamment abstrait pour répondre à toutes les problématiques ?
\end{block}
\begin{figure}
	\centering
	\includegraphics[height = 3cm]{img/ontologies}
\end{figure}
\end{frame}

\begin{frame}[c]{En guise d'exemple}
\vspace{-2em}
Travail sur les usages liés à l'arbre dans les socio écosystèmes
\begin{figure}
	\centering
	\includegraphics[height = 3cm]{img/ateliers}
  \includegraphics[height = 3cm]{img/bar_general}\\
  \includegraphics[height = 3cm]{img/class_CART}
\end{figure}
\end{frame}

%-=-=-=-=-=-=-=-=-=-=-=-=-=-=-=-=-=-=-=-=-=-=-=-=
%	FRAME: MERCI DE VOTRE ATTENTION
%-=-=-=-=-=-=-=-=-=-=-=-=-=-=-=-=-=-=-=-=-=-=-=-=
{
\usebackgroundtemplate{
\includegraphics[width=\paperwidth]{img/podor_fin}}%
\begin{frame}
  \vspace{-1em}
  \begin{minipage}[t][.8\textheight]{\textwidth}
    \color{\cnGrey}{\LARGE{Merci de votre attention !}}

    \vfill
  \end{minipage}
  \vspace{-3.5em}
  \centering
	You can find this presentation on github\includegraphics[height=0.85cm]{img/github}

\end{frame}
}


\end{document}
